\section{Related Works}
\label{sec:related_works}

\paragrapha{Visual Representations for MLLMs. }The visual encoder serves as a critical bridge between vision and language modalities in multimodal learning. Conventional multimodal language models typically employ CLIP-style models—such as CLIP~\cite{radford2021clip}, SigLIP~\cite{zhai2023siglip}, SigLIP2~\cite{tschannen2025siglip}, AIM~\cite{fini2025multimodal}, and DFN~\cite{fang2023dfn}—as visual encoders. However, such models often rely on fixed input resolutions, which constrains their flexibility in multimodal tasks, and their contrastively learned representations may not align optimally with downstream multimodal objectives. To address these limitations, several specialized MLLM visual encoders have been developed and trained end-to-end on multimodal tasks. These include open-source models like InternViT~\cite{chen2024internvl,zhu2025internvl3}, OryxViT~\cite{liu2024oryx}, and SAILViT~\cite{yin2025sail}, as well as proprietary encoders integrated into large MLLM systems such as Qwen-VL~\cite{Qwen2.5-VL}, Kimi-VL~\cite{team2025kimi}, and Seed-VL~\cite{guo2025seed1}. More recently, researchers have begun exploring unified discrete representations for vision and language. Quantized multimodal encoders such as QLIP~\cite{zhao2025qlip} and UniTok~\cite{ma2025unitok} apply quantization to image inputs. Nevertheless, these quantized visual encoders still exhibit a large performance gap compared to continuous models, particularly in tasks requiring textual understanding or fine-grained visual details. This gap can be attributed to the high compression ratio of discrete codes and the high sensitivity of multimodal models to detailed visual information.

\paragrapha{Quantized Visual Encoders. } Image tokenization is pivotal in bridging raw pixels with compact latent representations for visual generation. Among existing approaches, vector-quantized (VQ)~\cite{gersho2012vector} tokenizers have gained prominence due to their discrete latent space, which is well-suited for visual generation. The seminal VQ-VAE~\cite{van2017neural} introduced a learnable codebook to discretize continuous features via nearest-neighbor assignment. Subsequent methods (such as FSQ~\cite{mentzerfinite}, RPQ~\cite{chiu2022self}, BSQ~\cite{zhaoimage}, and LFQ~\cite{yulanguage}) explored structurally constrained codebooks with predefined geometries, often decoupled from end-to-end training. Enhancements like VQ-GAN~\cite{esser2021taming} further improved visual quality by incorporating adversarial and perceptual losses. While these tokenizers excel at learning compact, generative-friendly representations, they often struggle to preserve fine-grained visual details crucial for dense prediction and high-fidelity understanding tasks. To address this, we propose ViQ, a novel framework designed to minimize information loss and align discrete visual tokens with semantic and textual contexts.
